\title{Controlling Text-to-Image Diffusion by Orthogonal Finetuning}
For a block-diagonal orthogonal matrix divided into $r$ blocks, the number of learnable parameters becomes $\thickmuskip=2mu \medmuskip=2mu d(d/r-1)/2$, leading to a complexity of $\thickmuskip=2mu \medmuskip=2mu \mathcal{O}(d^2/r)$. The trend of adapting large pretrained models for downstream tasks has become increasingly prevalent~\cite{devlin2018bert,brown2020language,he2022masked}. \textbf{Why OFT converges faster}. Further clarification on these consistency metrics can be found in Appendix~\ref{app:settings}. To better understand this, consider approximating $\thickmuskip=2mu \medmuskip=2mu \bm{R}=(\bm{I}+\bm{Q})(I-\bm{Q})^{-1}$ with the Ne

As the series converges with respect to the operator norm, it is thus possible to embed the constraint $\thickmuskip=2mu \medmuskip=2mu\|\bm{R}-\bm{I}\|\leq\epsilon$ within the Cayley transform. Supplementary experimental results and methodological details are available in Appendix~\ref{app:settings}. Notably, OFT does not increase computational cost during inference. In addition to diffusion-based frameworks, alternatives such as GANs~\cite{reed2016generative,zhang2017stackgan,xu2018attngan,li2019controllable} and autoregressive models~\cite{ramesh2021zero,ding2021cogview,wu2022nuwa,yuscaling} have also been pursued for text-to-image tasks. Orthogonality can be achieved using the Cayley transform discussed in Section~\ref{sect:eff_ortho}. In contrast, ControlNet displays a more abrupt gain in controllability after the 8th epoch, a sudden convergence behavior previously noted by \cite{zhang2023adding}. Furthermore, OFT maintains computational efficiency by avoiding additional inference overhead, making it suitable for practical deployment. Results presented in Figure~\ref{fig:ortho} show that our implementation of OFT, when combined with orthogonality constraints, produces stable and accurately controlled results post training (at epoch 20). While LoRA shows superior performance over ControlNet on the S2I task, it underperforms on both the C2I and L2F tasks. While the latter is typically only implicit in standard finetuning routines, COFT makes this deviation constraint explicit. \section{Orthogonal Finetuning}

\subsection{Why Does Orthogonal Transformation Make Sense?}

To motivate orthogonal transformation in the context of finetuning text-to-image diffusion models, we divide the analysis into two aspects: (1) the rationale for adjusting neuron angles (that is, their directions) during finetuning, and (2) the reasons for choosing orthogonal transformations as the mechanism for these angular updates. During the training phase, we compute feature maps using the standard inner product, where $z$ is the output of the convolution kernel $\bm{w}$ applied to input $\bm{x}$ within a local window. Notably, no cosine-based activation (c) is involved in training—the model is trained exclusively with the inner product. Consequently, adopting a more structurally aware measure to compare the finetuned and pretrained models could substantially enhance both the safeguarding of generative capabilities and the stability of the finetuning process. For both DreamBooth and LoRA, implementations are based closely on their original descriptions, with hyperparameters set to the values recommended as optimal. Similarly, T2I-Adapter~\cite{mou2023t2i} introduces enhanced controllability by finetuning pretrained diffusion models. \textbf{Connection to LoRA}. This metric involves deriving the control signal from the generated image and then comparing it with the initial input control signal. Additionally, OFT and COFT deliver images with richer details and improved geometric coherence compared to ControlNet, all while utilizing significantly fewer parameters. Our experiments reveal that LoRA’s maximum attainable performance is quite limited, even after extensive hyperparameter optimization. Figure~\ref{fig:face_convergence} presents the trend of face landmark error as the models are finetuned over varying numbers of epochs. This enhanced stability in convergence makes OFT much more practical, as its training progression is notably more steady and foreseeable, facilitating the search for optimal hyperparameters. Through an extensive set of experiments, we observe substantial gains over existing approaches, including superior image quality, faster convergence, and greater finetuning stability. The forward computation in OFT can be reformulated as $\thickmuskip=2mu \medmuskip=2mu \bm{z}=(\bm{R}\bm{W}^0)^\top\bm{x}=(\bm{W}^0+(\bm{R}-\bm{I})\bm{W}^0)^\top\bm{x}$, where the term $\thickmuskip=2mu \medmuskip=2mu (\bm{R}-\bm{I})\bm{W}^0$ serves a similar role as LoRA’s low-rank adaptation. DreamBooth~\cite{ruiz2023dreambooth} advances this further by finetuning text-to-image diffusion models using a unique token for the subject and applying both reconstruction and class-specific prior losses, leveraging a limited set of subject images. The core principle involves learning an orthogonal transformation, shared across all neurons in a layer, so as to keep their inter-neuronal angles unaltered. This design preserves the original inference efficiency of the base pretrained model. \textbf{Quantitative comparison}. The value of the determinant of $\bm{R}$ ($1$ or $-1$) distinguishes a rotation from a reflection. \subsection{Controllable Generation}

\textbf{Settings}. The OFT framework gives rise to several compelling questions for future research. In the context of the C2I task, we utilize the IoU and F1 score as evaluation criteria. This qualitative assessment underscores that the OFT framework effectively balances both controllability and accurate subject representation. \textbf{Qualitative comparison}. Increasing $\epsilon$ grants more finetuning flexibility, making the COFT-adapted model behave more similarly to the OFT-based model. Similarly, \cite{nitzan2022mystyle} builds a personalized generative facial prior from multiple facial images of an individual. In the context of subject-driven generation, our protocol largely mirrors DreamBooth~\cite{ruiz2023dreambooth}, while for tasks involving controlled generation, we closely adhere to the procedures of ControlNet~\cite{zhang2023adding} and T2I-Adapter~\cite{mou2023t2i}. In line with the protocol in \cite{ruiz2023dreambooth}, we provide a quantitative assessment encompassing subject fidelity (utilizing DINO~\cite{caron2021emerging} and CLIP-I~\cite{radford2021learning}), alignment with textual prompts (CLIP-T~\cite{radford2021learning}), and diversity of the generated outputs (LPIPS~\cite{zhang2018unreasonable}). To limit this freedom, we advocate for the conservation of inter-neuronal angles, operating under the assumption that these angles encapsulate essential knowledge representation within neural networks. Additionally, we present a constrained version of this approach, which restricts angular shifts from the original pretrained weights to improve stability. Empirically, this restriction has negligible impact on practical performance. We also report results of a human evaluation in Appendix~\ref{app:human}, further substantiating OFT’s superior performance. Additional results, as well as expanded experimental discussions, are included in Appendix~\ref{app:settings},\ref{app:coft_oft},\ref{app:more_qualitative},\ref{app:failure}. However, integrating the $\epsilon$-deviation constraint within the Cayley-parametrized orthogonal transformation presents additional challenges. In the context of Transformers, LoRA is often integrated into the attention layers~\cite{hulora2022}. Additionally, the performance of OFT and COFT does not exhibit high sensitivity to the number of finetuning iterations, simplifying practical deployment by reducing the need for meticulous hyperparameter tuning across subjects. These findings substantiate that our OFT approach converges rapidly without sacrificing stability, specifically for subject-driven generation tasks. This metric measures how closely the control signal inferred from the generated output matches the original signal provided, further substantiating the effectiveness of OFT in enforcing controllability. Conversely, the OFT framework does not seem to have reached its performance limits, as improvements are still observed with increasing numbers of trainable parameters. Analogous to LoRA’s rank parameter $r'$, OFT introduces a block-diagonal parameter $r$ to limit the total number of parameters requiring optimization. This conservative choice of step size restricts how far the finetuned parameters can move from those of the pretrained network, effectively encouraging the solution to stay close to the original model by implicitly minimizing $\thickmuskip=2mu \medmuskip=2mu \| \bm{M} - \bm{M}^0\|$, where $\bm{M}$ and $\bm{M}^0$ represent the adapted and pretrained weights, respectively. For LoRA, which employs a rank-$r'$ decomposition, its count of additional trainable parameters is $\thickmuskip=2mu \medmuskip=2mu r'(d+n)$. CLIP-T reflects the average cosine similarity between the embeddings of text prompts and the generated images. \textbf{Quantitative comparison}. \textbf{Why not minimize hyperspherical energy}. Within the S2I task, LoRA fails to follow the guidance of the input segmentation map, whereas ControlNet, OFT, and COFT maintain accurate correspondence between the input and the generated outputs. Progressing to 2000 iterations, DreamBooth demonstrates image collapse, and LoRA erroneously produces yellow fur instead of a yellow shirt. For an orthogonal matrix of size $\thickmuskip=2mu \medmuskip=2mu d\times d$, the total number of free parameters is $\thickmuskip=2mu \medmuskip=2mu d(d-1)/2$, thus exhibiting a parameter complexity of $\mathcal{O}(d^2)$. \subsection{Re-scaled Orthogonal Finetuning}

We introduce a straightforward augmentation to the basic OFT framework by incorporating a learnable scaling parameter for each neuron. Stable Diffusion leverages VQ-GAN~\cite{esser2021taming} for its visual latent representation, while DALL-E2 utilizes CLIP~\cite{radford2021learning} for a multimodal latent space connecting images and text. We stress that our approach to quantitatively evaluating controllable generation constitutes a key contribution of our work, offering a robust measure of control fidelity in finetuned models, bounded by the accuracy of the pre-existing segmentation or detection models. The criteria for a robust finetuning approach can be outlined as follows: (1) \emph{training efficiency}, referring to minimal use of additional parameters and a reduced number of training epochs; and (2) \emph{generalizability preservation}, meaning the retention of both the fidelity and diversity of the model’s outputs. \textbf{Model finetuning}. \item We evaluate OFT on two core tasks: subject-driven generation and controllable generation. The block-diagonal parameterization, meanwhile, functions as an even more restrictive regularization for the orthogonal matrix. While this implicit regularization is reasonable, it may not suffice to constrain very large models. For the S2I task, we adopt mean IoU along with both mean and overall accuracy. To ensure a level playing field against LoRA, OFT or COFT is only applied to those network layers where LoRA is active. While Cayley parameterization significantly boosts efficiency, it introduces a minor constraint—only orthogonal matrices with determinant $1$, forming the special orthogonal group, can be represented. For the second question, an intuitive method to alter neuron directions involves applying a uniform rotation or reflection to all neurons within a layer, corresponding to an orthogonal transformation. Maintaining orthogonality of $\bm{R}$ can be accomplished using differentiable orthogonalization methods described in \cite{liu2021orthogonal,lezcano2019cheap}. In particular, we focus on two crucial finetuning scenarios for text-to-image generation: subject-driven generation, which seeks to produce subject-specific imagery using a small set of reference images and textual descriptions; and controllable generation, which augments the model’s input with explicit control signals. Their work establishes that orthogonal transformations possess ample expressivity to enable the training of robust classification models from scratch. Drawing on findings that hyperspherical similarity effectively encodes semantic content~\cite{liu2017deep,liu2018decoupled,chen2020angular}, we utilize hyperspherical energy~\cite{liu2018learning} to describe the pairwise relationships among neurons within a layer. \textbf{Finetuning stability and convergence}. All models under comparison are trained and evaluated using their respective optimal hyperparameters. Expanded results can be accessed in Appendix~\ref{app:more_qualitative}, \ref{app:more_control_task}, and \ref{app:failure}. Our principal contributions are summarized as follows:

\begin{itemize}[leftmargin=*,nosep]

    \item We introduce Orthogonal Finetuning, a new method designed to enhance controllability in text-to-image diffusion models during finetuning. Additionally, recognizing that reduced Euclidean distance between the finetuned and pretrained models tends to better conserve the benefits of pretraining, we present a variant named Constrained Orthogonal Finetuning~(COFT), which further restricts the adapted model to lie within a hypersphere of fixed radius centered on the original neurons. 
\begin{abstract}
Large-scale text-to-image diffusion models have demonstrated remarkable proficiency in synthesizing photorealistic images conditioned on textual prompts. In the C2I task, OFT and COFT manage to synthesize semantically coherent images from minimal Canny edge inputs, surpassing the performance of T2I-Adapter and LoRA. In contrast, OFT—though conceptually straightforward—effectively mediates this trade-off by maintaining the angles between neuron pairs, thereby achieving a favorable balance. However, these solutions frequently struggle to maintain fine subject-specific details. However, unlike LoRA, OFT applies a layer-shared orthogonal transformation, modifying neuron weights multiplicatively while ensuring that the relative angles between neurons within the same layer remain unchanged—thus offering greater stability. Additionally, OFT’s rapid convergence is corroborated by findings in Figure~\ref{fig:teaser}. These outcomes indicate that neuron angles (directions) are the principal carriers of semantic content from the input. To ensure this, the orthogonal matrix $\bm{R}$ is initialized as the identity matrix, aligning the initial fine-tuned weights with those of the pretrained model. Re-scaled OFT thus enhances the expressive capacity of OFT with an almost negligible increase in the number of trainable variables. We posit that an optimal finetuned model should retain a hyperspherical energy profile closely matching that of the pretrained model. Analysis of Figure~\ref{fig:dreambooth_final} indicates that both OFT and COFT maintain robust semantic subject fidelity, in contrast to LoRA, which frequently fails to retain subject identity (\emph{e.g.}, the subject is entirely unrecognizable in the bowl sample with LoRA). While traditional finetuning techniques allow maximum flexibility, they often lack robustness, frequently resulting in model instability or collapse. Moreover, it is possible to further compress the parameterization by sharing the block structure, that is, setting $\thickmuskip=2mu \medmuskip=2mu\bm{R}_i=\bm{R}_j,\forall i\neq j$. This is further supported by \cite{liu2021orthogonal}, demonstrating the adequacy of orthogonal transformations for neural network training. In line with these approaches~\cite{zhang2023adding,mou2023t2i}, our method incorporates OFT to finetune pretrained diffusion models, achieving superior controllability even with reduced training data and fewer tunable parameters. To further substantiate this perspective, we conducted an experiment evaluating the regularization effect brought by orthogonality constraints, following the experimental setup of ControlNet~\cite{zhang2023adding}. The example on the right in Figure~\ref{fig:teaser} illustrates OFT’s superior data efficiency relative to both ControlNet and LoRA, as OFT can achieve convergence with merely 5\% of the ADE20K dataset. \item Unlike prior finetuning strategies, OFT enables parameter updates while mathematically ensuring the conservation of hyperspherical energy—a property we empirically identify as key to maintaining the semantic fidelity of the original generative model. We further present a qualitative assessment of OFT and COFT, positioning them against leading techniques such as ControlNet, T2I-Adapter, and LoRA. To clarify this distinction, consider a weight matrix with dimensions $\thickmuskip=2mu \medmuskip=2mu 128\times 128$: with $\thickmuskip=2mu \medmuskip=2mu r'=8$, LoRA introduces $2,048$ trainable parameters; for $\thickmuskip=2mu \medmuskip=2mu r=8$ and no block sharing, OFT requires only $960$ trainable parameters. \section{Concluding Remarks and Open Problems}

Recognizing the significant role that angular relationships between neurons play in shaping visual semantics, we introduce a straightforward yet potent approach—orthogonal finetuning—for the manipulation of text-to-image diffusion models. For input $\thickmuskip=2mu \medmuskip=2mu \bm{x}\in\mathbb{R}^d$, the layer outputs $\thickmuskip=2mu \medmuskip=2mu \bm{z}=\bm{W}^\top\bm{x}\in\mathbb{R}^n$. Firstly, unlike \cite{liu2021orthogonal}, which actively minimizes hyperspherical energy, OFT is intended to preserve the hyperspherical energy observed in the pretrained model, thereby protecting its original architectural features during finetuning. The OFT framework is thus built upon this notion: that only orthogonal matrices, shared across each layer, are needed to finetune a neural network’s representations. Many current finetuning algorithms implicitly equate a reduced Euclidean distance between the finetuned and original models with superior retention of pretrained function, which often motivates the use of extremely small learning rates. Further conceptual grounding is provided by orthogonal over-parameterized training~\cite{liu2021orthogonal}, which initializes neural networks for classification by applying orthogonal transformations to randomly initialized weights. Furthermore, the OFT method exhibits resilience to the number of iterations, maintaining strong performance even with more iterations—a property not shared by DreamBooth or LoRA. Approaches such as \cite{avrahami2022blended,nichol2022glide} utilize user-provided masks to condition generation and protect subject regions. To illustrate, reconstructing an image by combining its original phase spectrum with a randomized magnitude spectrum can still yield the perceptual semantics of the original image. Diffusion models have more recently been leveraged for this purpose as well~\cite{saharia2022palette,voynov2022sketch,wang2022pretraining}. For our quantitative analysis, the evaluation metric is the mean $\ell_2$ distance measured between the control face landmarks and the landmarks predicted by each model. Among various finetuning strategies, adaptation techniques (\emph{e.g.}, \cite{hulora2022,houlsby2019parameter,pfeiffer2020adapterfusion}) have been extensively explored in natural language processing. For a fair comparison with LoRA, our implementation confines OFT to tuning the attention weights. This process can also be interpreted as redefining the canonical basis for the neurons in each layer. As illustrated by the samples in Figure~\ref{fig:control_final}, both OFT and COFT consistently generate images that are not only visually convincing and high-quality, but also exhibit precise controllability. This evidence affirms the essential role of the orthogonality constraint in the success of OFT. While OFT builds on the DreamBooth framework, it diverges by introducing orthogonal transformations for model finetuning instead of conventional small-step gradient updates. This adjustment is inspired by the insight that scaling neurons preserves hyperspherical energy, since normalization negates changes in magnitude. Guided by this property, we introduce Orthogonal Finetuning~(OFT), which allows large text-to-image diffusion models to adapt to new tasks while maintaining the original hyperspherical energy. Although block diagonal parametrization partially alleviates this, finding a more efficient, differentiable strategy for matrix inversion remains unresolved. Additional qualitative evidence can be found in Appendix~\ref{app:more_qualitative}. Our analysis indicates that maintaining this property is key to sustaining the semantic image generation capabilities inherent in text-to-image diffusion architectures. In comparison to previous techniques, OFT achieves higher controllability and more stable finetuning outcomes, while also requiring a smaller set of tunable parameters. For impartiality, image samples are selected at random across different approaches. Our method specifically addresses two major tasks: subject-driven generation and controllable generation. Inversion-based techniques~\cite{choi2021ilvr,dhariwal2021diffusion,ramesh2022hierarchical,gal2022image} offer the ability to change contextual elements while retaining the original subject. Conversely, reducing $\epsilon$ causes the COFT-tuned model to behave more like the original pretrained text-to-image diffusion model. This improved landscape for finetuning has also been observed empirically in \cite{liu2021orthogonal}. At test time, we assess three strategies for generating the feature map: (a) replicating the inner product used during training, (b) extracting only the magnitude, and (c) using purely angular information. The core challenge of finetuning is to strike an effective balance between adaptability and preserving model stability, and we contend that our OFT framework offers a principled solution to this trade-off. Investigating whether such compositions retain knowledge from all constituent downstream tasks presents an intriguing avenue for exploration. Figure~\ref{fig:toy_ae} reveals that reconstructing images from the angular components of neuron activations alone is nearly lossless, while magnitudes lack informativeness. In contrast with \cite{liu2021orthogonal}, our approach does not attempt to minimize hyperspherical energy. To assess the robustness and convergence rate during finetuning, we analyze DreamBooth, LoRA, OFT, and COFT, with the relevant results depicted in Figure~\ref{fig:teaser} and Figure~\ref{fig:db_convergence}. To ensure an equitable evaluation, every result is produced using the same finetuned model, with the respective method applied, and no additional optimization of the text prompt is performed for any final output. A central challenge remains: how to effectively steer or adapt these powerful models for a variety of downstream tasks. Motivated by this perspective, our approach involves learning an orthogonal transformation, shared across layers, that preserves the angular structure—thus maintaining constant hyperspherical energy for neurons within each layer. For reliability, each method undergoes finetuning on the identical pretrained model across 30 unique random seeds to diminish the influence of stochastic factors. Consistent with DreamBooth, all models are optimized using the same objective function. Furthermore, the training process of OFT can be interpreted as optimizing neuron weights constrained to the surface of a smooth hypersphere, which typically results in more favorable optimization landscapes. For classification tasks, having non-redundant neurons is preferable. Although this presumption may sometimes be valid, we contend that mere Euclidean closeness fails to provide a comprehensive picture of semantic preservation. Nevertheless, excessively modifying neuron directions without proper constraints tends to degrade the pretrained model’s generative abilities. Our observations indicate that, although OFT already performs strongly, introducing this explicit deviation constraint in COFT enhances finetuning stability beyond that of OFT. OFT can be viewed as seeking to keep the difference in hyperspherical energy between the finetuned and pretrained models minimal:

\begin{equation}\label{eq:oft_principle}
\footnotesize
\min_{\bm{W}} \left\| \text{HE}(\bm{W})-\text{HE}(\bm{W}^0)\right\|~~\Leftrightarrow~~\min_{\bm{W}} \bigg{\|} \sum_{i\neq j} \|\hat{\bm{w}}_i-\hat{\bm{w}}_j\|^{-1}-\sum_{i\neq j} \|\hat{\bm{w}}^0_i-\hat{\bm{w}}^0_j\|^{-1}\bigg{\|}
\end{equation}

Here, $\thickmuskip=2mu \medmuskip=2mu \hat{\bm{w}}_i:=\bm{w}_i/\|\bm{w}_i\|$ is the normalized vector for the $i$-th neuron, and the hyperspherical energy of $\bm{W}$ is defined as $\thickmuskip=2mu \medmuskip=2mu \text{HE}(\bm{W}):= \sum_{i\neq j} \|\hat{\bm{w}}_i-\hat{\bm{w}}_j\|^{-1}$. As depicted in Figure~\ref{fig:toy_ae}, the most direct way to affect semantic changes is by rotating neuron directions—a capability that OFT is designed to facilitate. We highlight two major distinctions between our Orthogonal Finetuning (OFT) method and the strategy from \cite{liu2021orthogonal}. Recent advances in text-to-image synthesis~\cite{saharia2022photorealistic,ramesh2022hierarchical,rombach2022high,gu2022vector,nichol2022glide} are largely attributable to the rapid evolution of diffusion models~\cite{ho2020denoising,song2021score,song2021denoising,dhariwal2021diffusion} and improved techniques in joint vision-language representations~\cite{su2020vlbert,lu2019vilbert,radford2021learning,wang2021simvlm,li2022blip,li2023blip,alayrac2022flamingo,schuhmann2022laion}. It is evident from these figures that both COFT and OFT maintain notable stability during diffusion model finetuning. Additionally, both OFT and COFT show significantly improved adherence to the conditioning text, whereas DreamBooth, despite reliably capturing subject identity, often struggles to generate outputs consistent with the prompt (\emph{e.g.}, in the bowl scenario, first row). Methods such as GLIDE~\cite{nichol2022glide} and Imagen~\cite{saharia2022photorealistic} operate within the pixel domain: GLIDE couples text encoder training with the diffusion process using text-image pairs, whereas Imagen incorporates a fixed, pretrained text encoder. The underlying justification is that random initialization yields a provably low expected hyperspherical energy, and the training methodology of \cite{liu2021orthogonal} is tailored to keep this energy low throughout learning (as lower energy is empirically linked to improved generalization in classification~\cite{liu2018learning,lin2020regularizing}). Through comprehensive experiments, we demonstrate that the OFT approach yields superior generation fidelity and more rapid convergence compared to prior methods. Furthermore, for COFT, an appropriate choice of $\epsilon$ ensures that both learning rate and iteration settings are largely straightforward. \end{abstract}

\section{Introduction}

The advent of modern text-to-image diffusion models~\cite{saharia2022photorealistic,ramesh2022hierarchical,rombach2022high} has set new benchmarks in text-driven, high-quality image creation. \end{itemize}

\section{Related Work}

\textbf{Text-to-image diffusion models}. Unlike methods such as ControlNet, our OFT framework imposes zero additional computational load at inference time on the finetuned model. Notably, regardless of these parameter reduction techniques, the resultant $\bm{R}$ is still orthogonal, ensuring the preservation of hyperspherical energy. Drawing inspiration from the zero initialization employed in LoRA, it is crucial for OFT to initiate fine-tuning from the pretrained parameters $\bm{W}^0$. Each method fine-tunes Stable Diffusion~(SD) v1.5 on all three datasets for 20 epochs. Secondly, OFT is explicitly designed to minimize deviations from the pretrained model, resulting in a constrained formulation; such constraints are absent in \cite{liu2021orthogonal}. Conversely, OFT and COFT continue to deliver reliable and coherent control of the outputs. Alternatively, if all blocks in $\bm{R}$ are shared, OFT amounts to transforming all neuron channels with a common orthogonal matrix. Figure~\ref{fig:eps_coft} illustrates the influence of the hyperparameter $\epsilon$ on COFT's performance. Finetuning approaches typically fall into two categories: directly modifying neuron weights using a low learning rate (\emph{e.g.}, \cite{ruiz2023dreambooth}), or integrating a compact module with reparameterized weights into the network (\emph{e.g.}, \cite{hulora2022,zhang2023adding}). Assuming both $r$ and $r'$ scale with the hidden size $d$ (for example, $\thickmuskip=2mu \medmuskip=2mu r=r'=\alpha d$ for some fixed $\thickmuskip=2mu \medmuskip=2mu 0<\alpha\leq 1$), LoRA's parameter count is then $\mathcal{O}(d^2+dn)$ while OFT's is $\mathcal{O}(d)$. Developing strategies to further enhance parameter efficiency in a manner that reduces such biases and improves adaptability is an open and significant challenge. The work by SphereNet~\cite{liu2017deep} further establishes that constraining neural network activations to a unit hypersphere not only retains representational capacity but can also improve generalization, suggesting that the orientation of neurons predominantly encodes crucial data attributes. By initializing $\bm{Q}$ as the zero matrix, the corresponding orthogonal matrix $\bm{R}$ is initialized as the identity, thereby providing identity initialization. From these results, it is clear that both COFT and OFT exhibit notably accelerated convergence. In contrast, approaches like Stable Diffusion~\cite{rombach2022high} and DALL-E2~\cite{ramesh2022hierarchical} train in the latent space. To provide a clearer insight into the advantages of OFT, we present a selection of randomly sampled results for subject-driven generation in Figure~\ref{fig:dreambooth_final}. In the special case where $\thickmuskip=2mu \medmuskip=2mu r=1$, the block-diagonal orthogonal structure degenerates into the conventional fully unconstrained orthogonal matrix. This results in an equivalent constraint of the form $\thickmuskip=2mu \medmuskip=2mu \|\bm{Q}-\bm{0}\|\leq\epsilon'$—where $\bm{0}$ is the zero matrix and $\epsilon'$ serves as a distinct error hyperparameter compared to $\epsilon$. To assess the effectiveness of controllable generation, we propose a control consistency metric. Image-to-image translation is fundamentally a controllable generation task, and historically, most solutions employ conditional generative adversarial networks~\cite{isola2017image,zhu2017toward,wang2018high,park2019semantic,choi2018stargan}. In particular, COFT enforces the following computation at inference:

\begin{equation}\label{eq:coft_general}
    \footnotesize
    \bm{z}=\bm{W}^\top\bm{x}=(\bm{R}\cdot\bm{W}^0)^\top\bm{x},~~~\text{s.t.}~\bm{R}^\top\bm{R}=\bm{R}\bm{R}^\top=\bm{I},~\norm{\bm{R}-\bm{I}}\leq \epsilon
\end{equation}

Here, the parameterization not only enforces orthogonality but also limits the deviation of $\bm{R}$ from the identity by at most $\epsilon$. \textbf{Subject-driven generation}. For improved efficiency, we utilize Cayley parameterization to construct the orthogonal matrix. Achieving minimum hyperspherical energy implies a uniform spacing of neurons over the hypersphere, but such an arrangement does not serve the objectives of finetuning since it risks erasing crucial information from pretraining. The experimental setup consists of three demanding controllable generation benchmarks presented in the main text: Canny edge to image~(C2I) on the COCO dataset~\cite{lin2014microsoft}, segmentation map to image~(S2I) on the ADE20K dataset~\cite{zhou2017scene}, and landmark to face~(L2F) on the CelebA-HQ dataset~\cite{karras2018progressive,xia2021tedigan}. DreamBooth~\cite{ruiz2023dreambooth} and LoRA~\cite{hulora2022} are chosen as our primary baselines. Unlike the original OFT forward mapping in Eq.~\eqref{eq:oft_general}, which adapts only $\bm{R}$, both the diagonal matrix $\bm{D}$ and the orthogonal matrix $\bm{R}$ are now subject to optimization. Secondly, OFT exhibits distinct promise for compositionality: the orthogonal matrices resulting from separate OFT finetuning processes can be composed via multiplication while preserving orthogonality. Our analysis reveals an inherent tension between flexibility and regularity during the finetuning process. \item \textbf{Controllable generation}~\cite{zhang2023adding,mou2023t2i}: This setting extends the generation process by incorporating additional guidance, such as canny edge maps or segmentation cues, alongside a textual description. We assess how rapidly ControlNet, T2I-Adapter, LoRA, and COFT converge on the L2F task, considering both quantitative and qualitative aspects. The work by \cite{hertz2022prompt} achieves local and global editing without requiring input masks. LoRA seeks to further restrict the parameter updates by imposing a low-rank condition: $\thickmuskip=2mu \medmuskip=2mu \text{rank}(\bm{M} - \bm{M}^0)=r'$, for a chosen small value $r'$. Importantly, OFT and COFT retain a comparable count of trainable parameters to LoRA (significantly less than ControlNet), yet they demonstrate much swifter convergence than existing alternatives. To illustrate the significance of neuron angles, we present a simplified experiment depicted in Figure~\ref{fig:toy_ae}, utilizing a conventional convolutional autoencoder trained on images of flowers. Nevertheless, parameterizing $\bm{R}$ using the Cayley transform can still be inefficient for high-dimensional cases (large $d$). In contrast, OFT maintains the orthogonality (and hence full-rank property) of the transformation applied to pretrained weights, thereby safeguarding the pretrained information from degradation. For instance, LoRA requires 20 epochs to achieve the accuracy that OFT attains after just 8 epochs. \end{document} \item To quantitatively assess controllable generation, we propose a control consistency metric. LPIPS is employed to determine the average pairwise cosine similarity for generated images, conditioned on the same subject and prompt. By contrast, removing the orthogonality constraint from OFT leads to failure in generating plausible images or achieving desired control. \subsection{Constrained Orthogonal Finetuning}

It is possible to further restrict the adaptability of the OFT approach by constraining the fine-tuned model to remain close to its initialization. In response to this limitation, this work focuses on two primary tasks within the realm of text-to-image generation:

\begin{itemize}[leftmargin=*,nosep]

    \item \textbf{Subject-driven generation}~\cite{ruiz2023dreambooth}: This task involves synthesizing images featuring a particular subject in novel contexts, given a handful of example images and a guiding text prompt. Within the initial 400 iterations, DreamBooth and OFT approaches successfully establish effective control over generation, in contrast to LoRA, which struggles to retain subject identity. Both methods allow the use of larger iteration counts without extensive manual tuning. \subsection{General Framework}

In standard finetuning protocols, model parameters (or selected layers) are updated using gradient descent with a limited learning rate. To enhance the robustness of the finetuning process, we additionally introduce Constrained Orthogonal Finetuning (COFT), which incorporates a radius constraint on the hypersphere. In contrast, our focus is on finetuning text-to-image diffusion models to enhance controllability and improve downstream generative capabilities. The outcomes demonstrate that our OFT framework consistently attains not only enhanced stability and convergence, but also achieves better performance across the considered evaluation metrics. \section{Experiments and Results}

\textbf{Experimental configuration}. This metric, which sums the hyperspherical similarity (such as cosine similarity) between every pair of neurons in a given layer, serves to measure neuron distribution uniformity across the unit hypersphere~\cite{liu2021learning}. The COFT method effectively integrates two explicit constraints: a minimal difference in hyperspherical energy and a bounded divergence from the pretrained parameters. OFT is most closely related to LoRA~\cite{hulora2022}, which proposes low-rank updates to model weights during finetuning. Although greater expressiveness can be reached via distinct angular adjustments for each neuron, our findings indicate that orthogonal transformation strikes a practical balance between model flexibility and regularization. Notably, OFT is a model-agnostic finetuning technique and can be integrated with any text-to-image diffusion architecture. In the L2F setting, we employ the average $\ell_2$ distance between the input control landmarks and those predicted. To mitigate this, we suggest expressing $\bm{R}$ as a block-diagonal matrix with $r$ constituent blocks, described by:

\begin{equation}
    \footnotesize
    \bm{R} = \text{diag}(\bm{R}_1, \bm{R}_2, \cdots, \bm{R}_r) = \begin{bmatrix}
    \bm{R}_1 \in \text{O}(\frac{d}{r}) &  &\\
    & \ddots & \\
    & & \bm{R}_r \in \text{O}(\frac{d}{r})
    \end{bmatrix} \in \text{O}(d)
\end{equation}

where $\text{O}(d)$ refers to the group of orthogonal transformations in $d$ dimensions, with $\thickmuskip=2mu \medmuskip=2mu \bm{R}\in\mathbb{R}^{d\times d}$ and each $\thickmuskip=2mu \medmuskip=2mu \bm{R}_i\in\mathbb{R}^{d/r\times d/r},\forall i$ being orthogonal matrices. We now compare OFT and LoRA with respect to parameter efficiency. \subsection{Subject-driven Generation}

\textbf{Experimental protocol}. While \cite{casanova2021instance} supports generating variations of specific instances, it may omit instance-defining features. In principle, OFT may be applied to any neural network architecture. \section{Intriguing Insights and Discussions}

\textbf{OFT is broadly applicable across architectures}. \textbf{No inference-time computational burden}. The introduction of ControlNet~\cite{zhang2023adding} marks a significant development, allowing additional control over pretrained diffusion models by finetuning them with supplementary control inputs, leading to strong performance in controlled generation. In particular, the orthogonal matrix is parameterized as $\thickmuskip=2mu \medmuskip=2mu \bm{R}=(\bm{I}+\bm{Q})(I-\bm{Q})^{-1}$, where $\bm{Q}$ is a skew-symmetric matrix that satisfies $\thickmuskip=2mu \medmuskip=2mu \bm{Q}=-\bm{Q}^\top$. In particular, the forward computation takes the form: $\thickmuskip=2mu \medmuskip=2mu\bm{z}=(\bm{R}\bm{W}^0\bm{D})^\top\bm{x}$,\footnote{\textbf{Errata}: The NeurIPS camera-ready submission incorrectly stated the forward pass for the re-scaled OFT as $\thickmuskip=2mu \medmuskip=2mu\bm{z}=(\bm{D}\bm{R}\bm{W}^0)^\top\bm{x}$. A fundamental obstacle is the absence of a metric to rigorously evaluate how well the pretrained generative prowess is retained. As comparative baselines, we employ ControlNet~\cite{zhang2023adding}, T2I-Adapter~\cite{mou2023t2i}, and LoRA~\cite{hulora2022}. Since $\bm{W}^0$ is usually of full rank, OFT also achieves a low-rank weight modification when $\thickmuskip=2mu \medmuskip=2mu \bm{R}-\bm{I}$ itself is low-rank. To address this, we exploit the invariance property of hyperspherical energy—specifically, that applying a uniform orthogonal transformation to all neurons guarantees preservation of their pairwise hyperspherical similarities. Lastly, the method’s parameter efficiency is fundamentally linked to its reliance on a block diagonal structure, which may introduce unwanted biases and curtail representational flexibility. While one straightforward strategy is to apply a regularization term to preserve hyperspherical energy throughout finetuning, such an approach does not reliably minimize the discrepancy in energy. LoRA employs a low-rank additive matrix to restrict significant alteration of information encoded in the pretrained weights. The DINO score follows a similar approach, but utilizes ViT S/16 DINO embeddings instead. \textbf{Controllable generation}. It is evident that adapter-based methods (such as T2I-Adapter and ControlNet) display slower convergence and achieve less satisfactory outcomes. Actively reducing this energy in diffusion model finetuning could lead to loss of semantic fidelity. Nevertheless, the semantic guidance provided by textual inputs may lack the specificity required for fine-level manipulation and detailed control of the resulting images. Stable Diffusion v1.5~\cite{rombach2022high} serves as the backbone text-to-image model in our evaluation. \end{itemize}

Both of these challenges center on the question of how to fine-tune text-to-image diffusion models in a way that maintains the original generative quality from pretraining. Projected gradient descent conveniently enforces the new constraint imposed on $\bm{Q}$. For further qualitative insights, we include extensive visual examples for each control task in Appendix~\ref{app:control_exp}, and also illustrate the versatility of OFT across additional control scenarios—such as dense pose to human body, sketch to image, and depth to image—in Appendix~\ref{app:more_control_task}. Notably, the phase spectrum encodes angular relationships between the input and its basis, which primarily retains the semantic content. Across all three control tasks, OFT and COFT consistently yield higher-quality images than their contemporary counterparts, highlighting the robustness and adaptability of the OFT framework for controllable generation. \subsection{Efficient Orthogonal Parameterization}\label{sect:eff_ortho}

Traditional orthogonalization approaches, like the differentiable Gram-Schmidt procedure, are known to be computationally intensive~\cite{liu2021orthogonal}. Qualitative examples in Figure~\ref{fig:face_convergence_qual} indicate that OFT and COFT both converge reliably, and the synthesized face poses progressively align with the intended control landmarks. \textbf{Convergence}. Firstly, OFT maintains orthogonality through Cayley parametrization, which necessitates the computation of a matrix inverse and slightly constrains the scalability of the approach. \textbf{Balancing flexibility and regularity in finetuning}. CLIP-I calculates the mean pairwise cosine similarity between CLIP embeddings of generated and real subject images. The statistics summarized in Table~\ref{table:dreambooth_final} indicate that both COFT and OFT deliver superior results over DreamBooth and LoRA on the DINO and CLIP-I metrics, showing significant improvements, and also slightly surpass or match these baselines in prompt fidelity and output diversity. Thus, altering semantic attributes of images is more likely achieved by adjusting neuron orientations rather than their magnitudes. We will elaborate on achieving this property efficiently. Moreover, there is insufficient theoretical grounding for constructing effective finetuning methods or for systematically selecting appropriate hyperparameters such as the number of epochs. The results presented in Table~\ref{table:control_gen} indicate that OFT and COFT surpass alternative methods in terms of both strength and accuracy of control. During inference, the learned orthogonal transformation $\bm{R}$ can be directly composed with the pretrained weight matrix $\bm{W}^0$ by simple matrix multiplication, yielding the updated parameter matrix $\thickmuskip=2mu \medmuskip=2mu \bm{W}=\bm{R}\bm{W}^0$. \textbf{Qualitative comparison}. For each approach, we report results from the model exhibiting the top validation CLIP scores. This block sharing reduces parameter complexity further to $\mathcal{O}(d^2/r^2)$. Specifically, setting the number of blocks equal to the number of input channels means each block operates on a single neuron channel, drawing an analogy to depth-wise convolution kernel learning~\cite{chollet2017xception}. Regarding the first question, our approach is motivated by prior empirical findings in \cite{liu2018decoupled,chen2020angular}, where it was observed that the angular difference between features effectively encapsulates the semantic gap. Nevertheless, the source code correctly implements the intended computation, so the results presented in Appendix~\ref{app:rescaled} remain accurate.} where $\thickmuskip=2mu \medmuskip=2mu \bm{D}=\text{diag}(s_1,\cdots,s_n)\in\mathbb{R}^{n\times n}$ is a diagonal matrix comprising positive learnable parameters $s_1,\ldots,s_n$. More formally, for a pretrained fully connected layer $\thickmuskip=2mu \medmuskip=2mu\bm{W}

\begin{equation}\label{eq:oft_general}
    \footnotesize
    \bm{z} = \bm{W}^\top \bm{x} = (\bm{R} \cdot \bm{W}^0)^\top \bm{x},~~~\text{s.t.}~\bm{R}^\top \bm{R} = \bm{R} \bm{R}^\top = \bm{I}
\end{equation}

In this formulation, $\bm{W}$ refers to the weight matrix fine-tuned using OFT, while $\bm{I}$ is the identity matrix. To address this issue, we propose a systematic finetuning technique called Orthogonal Finetuning~(OFT), aimed at customizing text-to-image diffusion models for specialized applications. Fundamentally, LoRA and OFT offer two fundamentally different parameter-efficient strategies: LoRA leverages low-rank decompositions to economize on parameter count, whereas OFT capitalizes on block-sparse orthogonal transformations to achieve similar goals. An illustration of OFT is provided in Figure~\ref{fig:block}. OFT also demonstrates improved sample efficiency by achieving strong performance with fewer training samples and epochs. It is straightforward to see that the minimal achievable value in Eq.~\eqref{eq:oft_principle} is zero, which occurs whenever $\bm{W}$ and $\bm{W}^0$ are related by an orthogonal transformation: $\thickmuskip=2mu \medmuskip=2mu \bm{W}=\bm{R}\bm{W}^0$ with $\thickmuskip=2mu \medmuskip=2mu \bm{R}\in\mathbb{R}^{d\times d}$ orthogonal. This observation implies that altering neuron orientations is fundamental to making semantic modifications in the generated output—a central aim in both subject-driven and controllable generative processes. Pivotal Tuning~\cite{roich2022pivotal} addresses this by finetuning the generator in the vicinity of an inverted latent representation, supplemented with regularization for identity preservation. We provide an initial empirical assessment of OFT applied to convolutional layer finetuning in Appendix~\ref{app:convolution}. For qualitative comparisons, our attention centers on OFT, COFT, and ControlNet, due to ControlNet's similar performance in terms of landmark error. To illustrate this, consider a dense layer $\thickmuskip=2mu \medmuskip=2mu \bm{W}=\{\bm{w}_1,\cdots,\bm{w}_n\}\in\mathbb{R}^{d\times n}$ where each column $\bm{w}_i\in\mathbb{R}^{d}$ represents a neuron, and $\bm{W}^0$ denotes its pretrained weights. Distinct from this, OFT enforces a constraint on the pairwise relationships among neurons, specifically by requiring $\thickmuskip=2mu \medmuskip=2mu \| \text{HE}(\bm{M})-\text{HE}(\bm{M}^0) \|=0$, where $\text{HE}(\cdot)$ denotes the hyperspherical energy corresponding to a weight matrix. The rationale behind utilizing orthogonal transformations for neuron finetuning draws partial analogy from the principles of the 2D Fourier transform, where an image can be separated into its magnitude and phase spectra. For experimental clarity, our main results utilize standard OFT to demonstrate the specific impact of orthogonal transformations, though our findings suggest re-scaled OFT generally achieves superior outcomes (see Appendix~\ref{app:rescaled}). OFT is also amenable to convolutional architectures, a scenario in which the block-diagonal structure of $\bm{R}$ offers compelling interpretations not found in LoRA. When evaluating the L2F task, our approaches deliver superior pose alignment in facial generation, particularly for difficult orientations. While straightforward, these strategies do not inherently ensure that the model’s pretraining generative capability is preserved. Distinct from previous approaches, OFT can theoretically guarantee the preservation of hyperspherical energy, an important measure of the pairwise relations among neurons mapped onto the unit hypersphere.